% 江苏海洋大学毕业设计（论文）开题报告（人文经管类）
% 对齐参考：工作手册 PDF 第 30-32 页，Word XML Table 13 + Table 14

\documentclass[a4paper,zihao=-4]{ctexart}

\RequirePackage{geometry}
\RequirePackage{../../styles/jouhandbook}

\geometry{
    a4paper,
    top=2.5cm,
    bottom=2.5cm,
    left=2.5cm,
    right=2.5cm
}

\begin{document}
\pagestyle{empty}

\begin{center}
    {\heiti\zihao{3}\textbf{毕业设计（论文）开题报告}}\\[0.2cm]
    {\heiti\zihao{-4}（人文经管类）}
\end{center}

\vspace{0.8cm}

{\zihao{-4}\songti
\begin{center}
\begin{tabular}{|L{3.33cm}|L{3.68cm}|L{2.23cm}|L{3.77cm}|}
\hline
题\hspace{1em}目: & \multicolumn{3}{L{9.67cm}|}{毕业设计（论文）题目（四号楷体\_GB2312下同）} \\
\hline
学\hspace{1em}院： & \multicolumn{3}{L{9.67cm}|}{商学院} \\
\hline
专业班级： & \multicolumn{3}{L{9.67cm}|}{国际经济与贸易\hspace{0.8em}国贸201} \\
\hline
学生姓名： & 扬帆起航 & 学号： & 2020000001 \\
\hline
指导教师： & \multicolumn{3}{L{9.67cm}|}{乘风破浪（职称）} \\
\hline
\end{tabular}
\end{center}
}

\vfill
\begin{center}
    \zihao{4}\songti 年\hspace{0.8cm}月\hspace{0.8cm}日
\end{center}

\newpage
\pagestyle{empty}

{\zihao{-4}\songti
1.课题研究的意义\par
\vspace{4.2cm}
2.课题的基本内容\par
\vspace{4.2cm}
3.课题的重点和难点\par
\vspace{4.2cm}
}

\newpage
\pagestyle{empty}

{\zihao{-4}\songti
\begin{center}
\begin{tabular}{|L{14.64cm}|}
\hline
4.论文提纲\par
\vspace{5.0cm} \\
\hline
指导教师意见（对课题的深度、广度及工作量的意见和对结果的预测）\par
\vspace{5.0cm}
\hfill 指导教师签字：\hspace{3.5cm}年\hspace{0.7cm}月\hspace{0.7cm}日 \\
\hline
系意见：\par
\vspace{4.4cm}
\hfill 系主任签字：\hspace{3.5cm}年\hspace{0.7cm}月\hspace{0.7cm}日 \\
\hline
\end{tabular}
\end{center}
}

\end{document}
